\chapter{\IfLanguageName{dutch}{Control-flow visualisatie}{Control-flow visualisation}}%
\label{ch:cfvisualisatie}

Dit hoofdstuk beschrijft hoe de control-flow van HLASM-programma’s gevisualiseerd kan worden. Het doel van deze visualisatie is om beginnende ontwikkelaars een overzicht te geven van de mogelijke control-flow binnen een HLASM-programma.

In HLASM is de programmastructuur niet onmiddellijk zichtbaar in de broncode. De control-flow moet daarom afgeleid worden op basis van labels, spronginstructies en de sequentiële uitvoering van instructies. Een grafische representatie maakt deze structuur beter interpreteerbaar.

De visualisatie wordt gerealiseerd in de vorm van een webapplicatie. Dit maakt de analyse en visualisatie laagdrempelig en interactief, zonder bijkomende installatievereisten. De gebruiker kan HLASM-code invoeren in de browser, waarna automatisch een control-flowdiagram wordt gegenereerd.

\section{Requirements}

Voor het ontwikkelen van een proof-of-concept worden zowel functionele als niet-functionele vereisten opgesteld. Deze vereisten bepalen welke onderdelen van HLASM-code worden geanalyseerd en hoe de visualisatie wordt opgebouwd.

\subsection{Functionele vereisten}

Voor de PoC zijn de volgende functionele vereisten gedefinieerd:

\begin{itemize}
    \item Detectie van labels in de broncode, zodat sprongdoelen geïdentificeerd kunnen worden.
    \item Detectie van vertakkingsinstructies die de control-flow beïnvloeden, waarondedr zowel conditionele als onvoorwaardelijke sprongen.
    \item Correcte mapping van labels naar knooppunten (nodes) in de control-flowgrafiek.
    \item Correcte mapping van sprongen naar verbindingen (edges) tussen knooppunten.
    \item Ondersteuning van meerdere uitvoeringspaden binnen één programma, inclusief eenvoudige vertakkingen en lussen.
\end{itemize}

\subsection{Niet-functionele vereisten}

Naast de functionele vereisten zijn er ook niet-functionele vereisten opgesteld:

\begin{itemize}
    \item Leesbaarheid: de gegenereerde diagrammen moeten een duidelijk overzicht geven van de control-flow, waarbij vertakkingen en uitvoeringspaden eenvoudig te volgen zijn.
    
    \item Correctheid: de gegenereerde Mermaid.js-code moet syntactisch correct zijn en zonder aanpassingen renderbaar zijn.
    
    \item Automatisatie: de visualisatie moet volledig automatisch gegenereerd worden op basis van de HLASM-broncode, zonder manuele tussenstappen.
    
    \item Toegankelijkheid: de toepassing wordt aangeboden als webapplicatie, zodat deze onmiddellijk bruikbaar is via een browser zonder installatie.
    
    \item Beperkingen: de implementatie richt zich op een afgebakende subset van HLASM-instructies. 
    
\end{itemize}


\section{Ontwerp}

Het ontwerp van de visualisatietool bestaat uit een vertaling van HLASM-broncode naar een grafische representatie van de control-flow. Deze representatie wordt getoond als een flowchart.

De implementatie gebeurt als een webapplicatie. De logica wordt geimplementeerd via JavaScript. Hierdoor blijft de toepassing eenvoudig en onafhankelijk van externe systemen.

De opbouw van de flowchart gebeurt in verschillende stappen:

\begin{itemize}
    \item De ingevoerde broncode wordt opgesplitst in individuele instructies.
    \item Per instructie worden het label, de mnemonic en de operanden geïdentificeerd op basis van de kolomgebaseerde structuur.
    \item Alle labels worden verzameld en gekoppeld aan hun positie in de code.
    \item Voor elke instructie wordt bepaald naar welke instructie de uitvoering kan verdergaan.
\end{itemize}

\newpage

Hierbij wordt onderscheid gemaakt tussen expliciete en impliciete control-flow:

\begin{itemize}
    \item \textbf{Expliciete control-flow}: veroorzaakt door branch-instructies die naar een label springen.
    \item \textbf{Impliciete control-flow}: de standaard sequentiële uitvoering naar de volgende instructie.
\end{itemize}

\subsection{Mapping van HLASM naar Mermaid.js}

Om de control-flowgrafiek visueel weer te geven, wordt gebruik gemaakt van Mermaid.js. Hiervoor wordt een mapping gedefinieerd tussen HLASM-elementen en grafische elementen.

\begin{table}[h!]
    \centering
    \small
    \begin{tabular}{|p{3.5cm}|p{3.5cm}|p{5cm}|}
        \hline
        \textbf{HLASM-element} & \textbf{Mnemonic} & \textbf{Mermaid.js representatie} \\
        \hline
        Label & Niet van toepassing & Node in de flowchart \\
        \hline
        Sequentiële instructie & L, AR, SR, etc. & Verbinding naar volgende node \\
        \hline
        Onvoorwaardelijke sprong & B, BR & Edge naar doelnode \\
        \hline
        Conditionele sprong & BE, BNE, BH, BL & Vertakkende edge met conditie \\
        \hline
        Programma start & CSECT, START & Entry node \\
        \hline
        Programma einde & END & Eindnode \\
        \hline
    \end{tabular}
    \caption{Mapping van HLASM-elementen naar Mermaid.js}
    \label{tab:mapping_hlasm_mermaid}
\end{table}

Bij conditionele sprongen ontstaan meerdere mogelijke paden. In de visualisatie worden deze weergegeven als vertakkingen met een label dat de voorwaarde aanduidt. Onvoorwaardelijke sprongen worden voorgesteld als directe verbindingen zonder extra voorwaarden.

Sequentiële instructies worden impliciet verbonden met de volgende instructie in de code, tenzij deze onderbroken worden door een spronginstructie.

Door deze mapping toe te passen kan de structuur van een HLASM-programma worden weergegeven als een flowchart, waarbij de nadruk ligt op de control-flow.

\newpage

\subsubsection{Voorbeeld}

Om de mapping tussen HLASM en Mermaid.js te verduidelijken, wordt een eenvoudig voorbeeld gegeven.

\begin{lstlisting}[language=HLASM, caption={Voorbeeld van een conditionele vertakking.}]
    ----+----1----+----2----+----3----+----4----+----5----+----6----+----7--
             CR    R1,R2
             BE    GELIJK
             B     VERDER
    GELIJK   ...
    VERDER   ...
\end{lstlisting}

Dit voorbeeld kan worden omgezet naar volgende Mermaid.js representatie:

\begin{lstlisting}[language={}, caption={Gegenereerde Mermaid.js flowchart.}]
    flowchart TD
    A[CR R1,R2] -->|BE| B[GELIJK]
    A --> C[B VERDER]
    C --> D[VERDER]
    B --> D
\end{lstlisting}
